% Generated draft: requires mathematical and source review.
\documentclass[11pt]{article}
\usepackage[margin=25mm]{geometry}
\usepackage{fontspec}
\setmainfont{Latin Modern Roman}
\usepackage{amsmath,amssymb,mathtools,unicode-math}
\setmathfont{Latin Modern Math}
\usepackage{enumitem,needspace,xurl,hyperref}
\setlength{\parindent}{0pt}
\setlength{\parskip}{5pt}
\setlength{\emergencystretch}{3em}
\newcommand{\N}{\mathbb{N}}
\newcommand{\Z}{\mathbb{Z}}
\newcommand{\Q}{\mathbb{Q}}
\newcommand{\R}{\mathbb{R}}
\newcommand{\C}{\mathbb{C}}
\newcommand{\F}{\mathbb{F}}
\newcommand{\A}{\mathbb{A}}
\newcommand{\PP}{\mathbb{P}}
\newcommand{\E}{\mathbb{E}}
\newcommand{\eps}{\varepsilon}
\newcommand{\dd}{\,\mathrm{d}}
\begin{document}
\section*{UnsolvedMath ID: 701415}
\textbf{Source problem code:} catalogue entry
\subsection{Definitions and mathematical statement}
Let \(A\) and \(B\) be nontrivial groups, let \(C\) be a group, and let \(\iota_A:C\hookrightarrow A\) and \(\iota_B:C\hookrightarrow B\) be injective homomorphisms. Their free product with amalgamation is
\[
G=A\mathbin{*_{C}}B
=(A*B)\big/\left\langle\!\left\langle
\iota_A(c)\iota_B(c)^{-1}:c\in C
\right\rangle\!\right\rangle,
\]
where \(A*B\) is the free product and the double brackets denote normal closure in \(A*B\). A maximal subgroup of \(G\) means a proper subgroup \(M<G\) such that, for every subgroup \(H\leq G\),
\[
M\leq H\leq G\quad\Longrightarrow\quad H=M\ \text{or}\ H=G.
\]
The question attributed to J. Wiegold asks whether every generalized free product with amalgamation of two nontrivial groups has at least one maximal subgroup. Under the unrestricted interpretation of amalgamation, this asks whether
\[
\forall A,B,C,\iota_A,\iota_B\text{ as above},\qquad
\exists M<G\quad M\text{ is a maximal subgroup of }G.
\]
The supplied wording does not specify whether the amalgamated subgroup must be proper in both factors. Consequently, the displayed unrestricted quantification is a conditional interpretation of the source: if proper amalgamation is intended, the hypotheses must additionally require \(\iota_A(C)\neq A\) and \(\iota_B(C)\neq B\). The intended convention requires review. No finiteness or countability assumption is supplied.
\subsection{Short English statement}
Does every free product with amalgamation of two nontrivial groups contain a maximal proper subgroup? The source leaves unclear whether the amalgamated subgroup must be proper in both factors.
\subsection{Sources}
\begin{itemize}
\item Local UnsolvedMath registry, permanent ID \(701415\), exact source code \texttt{catalogue entry}, title \emph{Kourovka Notebook Problem 17.114}; supplied upstream revision \texttt{d371ee917ac9647fbb06cd6fc529deec6f3dda82}.
\item \emph{The Kourovka Notebook (Unsolved Problems in Group Theory)}, original problem \(17.114\), proposed by J. Wiegold, 17th issue (2010). The supplied background also identifies the 21st issue edition (2026), \texttt{arXiv:1401.0300v48}. The supplied collection description names E. I. Khukhro and V. D. Mazurov as editors.
\item The local record labels the problem \emph{open}. This is supplied source metadata, not an independently verified assessment of its mathematical or literature status. The cited source has not been independently consulted for this draft.
\end{itemize}
\end{document}
